\section{Conclusion}\label{sec:conclusion}

We presented ContamCal, a framework that connects contamination detection signals to IRT-based calibration through a two-stage pipeline: training-free signal fusion via RRF followed by contamination-augmented IRT estimation. Our experiments on MMLU, ARC-Challenge, and GSM8K with controlled contamination injection across the Qwen2.5 family yield three findings. First, IRT ability estimation provides calibration benefit independent of the contamination parameter $\gamma$, reducing calibration error by up to 53\% relative to raw scores. This decomposition shows that the primary value of psychometric modeling for contaminated benchmarks lies in separating stable ability from item-level noise. Second, Min-K\%++ exhibits directional reversal at small model scales, producing actively misleading conclusions rather than merely failing to detect contamination. Third, simulation studies and empirical results at $J = 13$ confirm that reliable per-domain $\gamma$ recovery requires ensemble sizes ($J \geq 50$) well above current LLM evaluation practice; exploratory domain analysis reveals directional heterogeneity across MMLU cognitive domains, but most estimates are underpowered at current sample sizes.

The central lesson is that IRT provides a principled calibration framework for contaminated benchmarks, with per-domain $\gamma$ parameters offering diagnostic potential that will become increasingly actionable as evaluation pools grow. Scaling to $J \geq 50$ through broader model sampling, multidimensional IRT for domains where scalar $\gamma$ is insufficient, and active signal selection for closed-source settings are immediate next steps.